\documentclass[10pt]{amsart}
\usepackage[a4paper,margin=23mm]{geometry}
\usepackage{amsmath,amssymb,amsthm,xfrac,hyperref,xspace,graphicx}
\usepackage[english]{babel}
\usepackage[scr=rsfs,cal=euler]{mathalfa}
\newcommand{\resp}{resp.}
\newcommand{\cf}{cf.}
\newcommand{\ie}{i.e.}
\newcommand{\eg}{e.g.}
\newcommand{\loccit}{loc.\ cit.}
\newcommand{\skpt}{\par\smallskip\noindent}
\let\Subsection\subsection
\let\Subsubsection\subsubsection
\emergencystretch=3em
\pagestyle{plain}
\graphicspath{{hussain-lamzouri_figures}}

\multlinegap0pt



\theoremstyle{plain}
\newtheorem{thm}{Theorem}[section]
\newtheorem{cor}[thm]{Corollary}
\newtheorem{lem}[thm]{Lemma}
\newtheorem{pro}[thm]{Proposition}

\theoremstyle{definition}
\newtheorem{rem}[thm]{Remark}

\newcommand{\Legendre}[2]{\Bigl(\dfrac{#1}{#2}\Bigr)}
\newcommand{\legendre}[2]{\bigl(\tfrac{#1}{#2}\bigr)}

\let\ep\varepsilon

\newcommand{\X}{\mathbb{X}}
\newcommand{\Y}{\mathbb{Y}}
\newcommand{\ex}{\mathbb{E}}
\newcommand{\pr}{\mathbb{P}}
\newcommand{\B}{\mathcal{B}}
\newcommand{\un}{\boldsymbol{n}}
\newcommand{\ut}{\boldsymbol{t}}

\begin{document}
\title{Unconditional joint-moment asymptotics for Legendre paths averaged over dyadic prime intervals}
\author{Ayesha Hussain \and Youness Lamzouri}
\maketitle
\section{Source definitions}
\begin{equation}\label{eq: definition of character sum}
S_p(x) := \sum_{n \leq x} \Legendre{n}{p},
\end{equation}
where $\bigl(\frac{\cdot}{p}\bigr)$ is the Legendre symbol modulo an odd prime $p$ and $1\leq x\leq p$.
\begin{equation}\label{Eq: RelationFpSp}
f_p(t) = \frac{1}{\sqrt{p}}S_p(tp) + \frac{\{ pt \}}{\sqrt{p}} \Legendre{\lceil pt \rceil }{p}= \frac{1}{\sqrt{p}}S_p(tp) + O\Bigl(\frac{1}{\sqrt{p}}\Bigr),
\end{equation}
where $\{x \}$ is the fractional part of $x$.

In this paper, we shall prove that these paths, viewed as random processes on $[0, 1]$, have a nice limiting distribution. Before stating our result, we shall construct a probabilistic random model for the Legendre path $f_p$.
A good random model for the Legendre symbol, which was extensively studied in recent years, is a \emph{Rademacher random completely multiplicative function}. This is defined as\vspace*{-3pt}
\begin{equation}\label{eq: RademacherRandom}\X_n:= \prod_{q^k|| n} \X_q^k
\end{equation}
for $n\in \mathbb{N}$, where $\X_1=1$, and $\{\X_q\}_{q \text{ prime}}$ is a sequence of I. I. D. random variables taking the values $\pm 1$ with probability $1/2$ each. The distribution of sums of Rademacher random multiplicative functions was extensively studied in recent years, in particular by Harper \cite{Ha1} and \cite{Ha2}. In studying the distribution of the character sum $S_p(x)$, our first guess is to model the normalized sum $\frac{1}{\sqrt{p}}S_p(tp)$ by the following normalized sum of random variables
\[
\frac{1}{\sqrt{p}}\sum_{n\leq tp }\X_n.\]
However, it turns out that this is not a good model for $\frac{1}{\sqrt{p}}S_p(tp)$, since it does not account for the periodicity of the Legendre symbol modulo $p$. To exploit this periodicity, we shall use the Fourier series expansion of $f_p$, which was first established by Pólya in the following quantitative form \cite[eq.\,(9.19), p.\,311]{MoVaBook}
\begin{equation}\label{eq: definition Legendre path}
f_p(t) = \frac{\tau(\legendre{\cdot}{p})}{2 \pi i \sqrt{p}} \sum_{1 \leq |n| \leq Z} \Legendre{n}{p} \frac{1 - e(-nt)}{n} + O \Bigl(\frac{1}{\sqrt{p}}+\frac{\sqrt{p}\log p}{Z} \Bigr),
\end{equation}
for any real number $Z\geq 1$, where $e(t):= e^{2\pi i t}$ and $\tau(\Legendre{\cdot}{p})$ is the Gauss sum associated to the Legendre symbol modulo $p$, more precisely
\begin{equation}\label{eq: GaussSum}
\begin{aligned}
\tau\Bigl(\Legendre{\cdot}{p}\Bigr) := \sum_{a = 1}^p \Legendre{a}{p} e(a/p) = \begin{cases}
\sqrt{p}, & \text{if } p \equiv 1 \pmod 4, \\
i \sqrt{p}, & \text{if } p \equiv 3 \pmod 4.
\end{cases}
\end{aligned}
\end{equation}
Inserting the values of the Gauss sum \eqref{eq: GaussSum} in the Fourier expansion \eqref{eq: definition Legendre path} gives
\begin{equation}\label{eq: PolyaFourier}
f_p(t) = \frac{\varepsilon_p}{2 \pi i } \sum_{1 \leq |n| \leq Z} \Legendre{n}{p} \frac{1 - e(-nt)}{n} + O \Bigl(\frac{1}{\sqrt{p}}+\frac{\sqrt{p}\log p}{Z} \Bigr),
\end{equation}
where $\varepsilon_p= 1$ if $p\equiv 1\pmod 4$, and equals $i$ if $p\equiv 3\pmod 4$. Let $\{\X_q\}_{q \text{ prime}}$ be a sequence of I. I. D. random variables taking the values $\pm 1$ with equal probability $1/2$. Since the primes split evenly into the residue classes $1$ and $3 \pmod 4$, we shall model the value $\bigl(\frac{-1}{p}\bigr)$ by a random variable $\X_{-1}$, which is independent from the $\X_q$'s for $q$ prime, and takes the values $\pm1$ with equal probability $1/2$. We also extend the definition of Rademacher random completely multiplicative functions to the negative integers using multiplicativity, namely by setting $\X_{-n}=\X_{-1}\X_n$ for $n\in \mathbb{N}$. Combining all these elements together leads us to model the Legendre path $f_p$ by the following random Fourier series\vspace*{-3pt}
\[
F_{\X}(t):= \frac{\Y}{2 \pi i} \sum_{n \in \mathbb{Z}\setminus\{0\}} \X_n \frac{1 - e(-nt)}{n} \]
for $t\in [0,1]$, where the random variable $\Y$ is defined by $\Y= 1$ if $\X_{-1}=1$, and $\Y=i$ if $\X_{-1}=-1$. For a fixed $t\in [0, 1]$, it follows from \cite[Lem.\,1]{Ka} that the series defining $F_{\X}(t)$ converges almost surely as the limit of the partial sums\vspace*{-3pt}
\[F_{\X, N}(t):=\frac{\Y}{2 \pi i} \sum_{0<|n|\leq N} \X_n \frac{1 - e(-nt)}{n}. \]
and $\pi^*(Q)$ is the number of primes $p$ in the interval $[Q, 2Q]$.
\begin{rem} Recall that the behavior of the Legendre path attached to $p$ depends on the value of the Legendre symbol $\legendre{-1}{p}$. Indeed, it follows from \eqref{eq: PolyaFourier} that for any real number $Z\geq 1$ we have
\begin{equation} \label{eq: Polya1mod4}f_p(t) = \frac{1}{\pi } \sum_{1 \leq |n| \leq Z} \Legendre{n}{p} \frac{\sin(2\pi n t)}{n} + O\Bigl(\frac{1}{\sqrt{p}}+\frac{\sqrt{p}\log p}{Z}\Bigr),
\end{equation}
if $p\equiv 1 \pmod 4$, and
\begin{equation} \label{eq: Polya3mod4}f_p(t) = \frac{1}{\pi } \sum_{1 \leq |n| \leq Z} \Legendre{n}{p} \frac{1-\cos(2\pi n t)}{n} + O\Bigl(\frac{1}{\sqrt{p}}+\frac{\sqrt{p}\log p}{Z}\Bigr),
\end{equation}
if $p\equiv 3\pmod 4$. As before, we can view each $p \mapsto f_p$ for $Q\leq p\leq 2Q$ and $p\equiv 1\pmod 4$ (respectively $p\equiv 3\pmod 4$) as a random variable on the finite probability space $\{Q\leq p\leq 2Q \text{ and } p\equiv 1\pmod 4\}$ (respectively $\{Q\leq p\leq 2Q \text{ and } p\equiv 3\pmod 4\}$) endowed with the uniform probability measure, and we denote this random variable by $\mathcal{F}_{Q, +}$ (respectively $\mathcal{F}_{Q, -}$). We also
do the same with the random Fourier series, where we fix $\X_{-1} = 1$ or $-1$. To this end we define for $t\in [0, 1]$
\begin{align*}
&F_{\X,+}(t) = \frac{1}{\pi} \sum_{n \geq 1} \X_n \frac{\sin(2\pi nt)}{n}, &\quad &F_{\X,-}(t) = \frac{1}{\pi} \sum_{n \geq 1} \X_n \frac{1 - \cos(2 \pi n t)}{n}.
\end{align*}
Then, it follows from the proof of Theorem \href{https://doi.org/10.5802/jep.260}{1.1 of the source article} that the sequence of processes $(\mathcal{F}_{Q, \pm})_{Q}$ converges to the process $F_{\X, \pm}$.
\end{rem}
\section{Joint moment definitions}
\setcounter{subsection}{1}
\Subsection{The sequence of joint moments of the random process is determinate}\label{s: the limiting moment}
In Section \ref{s: proof of convergence in finite distributions real} we shall use the method of moments to prove that the sequence of processes $(\mathcal{F}_{Q})_Q$ converges to the process $F_{\X}$ in the sense of convergence of finite distributions, as $Q\to \infty$. More precisely, given an integer $k \geq 1$ and a $k-$tuple $0 \leq t_1 < \cdots < t_k \leq 1$, we will prove that the joint moments of $(\mathcal{F}_{Q}(t_1), \dots, \mathcal{F}_{Q}(t_k)) $
converge to those of $(F_{\X}(t_1), \dots, F_{\X}(t_k))$ as $Q\to \infty$. However, we first need to verify that these moments are determinate, i.e. that they have only one representing measure. For $\boldsymbol{n}=(n_1, \dots, n_k)\in (\mathbb{Z}_{\geq 0})^k$ we let $n= \sum_{j=1}^k n_j$, and define
\begin{equation}\label{eq: limiting moment}
M_{\X}(\boldsymbol{n}):= \mathbb{E}\biggl( \prod_{i = 1}^k F_{\X}(t_i)^{n_i} \biggr), \quad
M_{\X,\pm}(\boldsymbol{n}):= \mathbb{E}\biggl( \prod_{i = 1}^k F_{\X,\pm}(t_i)^{n_i} \biggr).
\end{equation}
Then we have
\begin{equation}\label{eq: MomentRandomPM}
\begin{aligned}
M_{\X}(\boldsymbol{n})&= \mathbb{E}\biggl( \prod_{i = 1}^k F_{\X}(t_i)^{n_i} \ \Big| \ \X_{-1}=1\biggr) \pr(\X_{-1}=1) \\
& \quad \quad \quad +\mathbb{E}\biggl( \prod_{i = 1}^k F_{\X}(t_i)^{n_i} \ \Big| \ \X_{-1}=-1\biggr) \pr(\X_{-1}=-1)\\
& =\frac{1}{2}\left(M_{\X,+}(\boldsymbol{n})+ M_{\X,-}(\boldsymbol{n})\right),
\end{aligned}
\end{equation}
since $\pr(\X_{-1}=1)= \pr(\X_{-1}=-1)=1/2.$
Therefore, using the definition of $F_{\X,+}(t)$ and expanding the moment $M_{\X,+}(\boldsymbol{n})$ we obtain
\begin{equation}\label{eq: FormulaMomentRandom+}
\begin{aligned}
M_{\X,+}(\boldsymbol{n}) &= \frac{1}{\pi^{n}} \mathbb{E} \biggl( \sum_{a_{1,1}, \dots, a_{k,n_k} \geq 1} \prod_{i = 1}^k \prod_{j = 1}^{n_i} \X_{a_{i,j}} \frac{\sin(2 \pi a_{i,j}t_i)}{a_{i,j}}\biggr)\\ & = \frac{1}{\pi^{n}} \mathbb{E} \biggl(\sum_{a=1}^{\infty} \frac{\X_a}{a}\mathcal{B}_{\boldsymbol{n},\mathbf{t}, +}(a)\biggr)= \frac{1}{\pi^{n}} \sum_{a=1}^{\infty} \frac{\mathcal{B}_{\boldsymbol{n},\mathbf{t}, +}(a^2)}{a^2},
\end{aligned}
\end{equation}
where
\begin{align*}
\mathcal{B}_{\boldsymbol{n},\mathbf{t},+}(a)
:= \sum_{a_{1}\cdots a_{k} = a} \prod_{i = 1}^k \sum_{b_{i,1} \cdots b_{i,n_i} = a_i } \prod_{j = 1}^{n_i} \sin(2 \pi b_{i,j} t_i).
\end{align*}
We observe that uniformly for $\mathbf{t}$ we have
\begin{equation}\label{eq: BoundBnt+}
|\mathcal{B}_{\boldsymbol{n},\mathbf{t}, +}(a)| \leq \sum_{a_1 \cdots a_k = a} \prod_{i = 1}^k d_{n_i}(a_i) = d_n (a),
\end{equation}
where $d_n$ is the $n$-th divisor function.
Similarly, one has
\begin{equation}\label{eq: FormulaMomentRandom-}
M_{\X,-}(\boldsymbol{n})= \frac{1}{\pi^{n}} \sum_{a=1}^{\infty} \frac{\mathcal{B}_{\boldsymbol{n},\mathbf{t}, -}(a^2)}{a^2},
\end{equation}
where \begin{align*}
\mathcal{B}_{\boldsymbol{n},\mathbf{t},-}(a)
:= \sum_{a_{1}\cdots a_{k} = a} \prod_{i = 1}^k \sum_{b_{i,1}\cdots b_{i,n_i} = a_i } \prod_{j = 1}^{n_i} (1-\cos(2 \pi b_{i,j} t_i)).
\end{align*}
We also have the following analogous bound to \eqref{eq: BoundBnt+}
\begin{equation}\label{eq: BoundBnt-}|\mathcal{B}_{\boldsymbol{n},\mathbf{t}, -}(a)|\leq 2^nd_n(a).
\end{equation}
\section{Unconditional joint moments}\label{s: proof of convergence in finite distributions real}
\setcounter{thm}{1}
\begin{pro}\label{thm: moments theorem}
Let $A>1$ be a fixed constant. Let $k$ be a positive integer, $\boldsymbol{n} = (n_1, \dots, n_k)\in (\mathbb{Z}_{\geq 0})^k$ and $0 \leq t_1 < \cdots < t_k \leq 1$ be real numbers. Then we have
\begin{equation}\label{AsymptoticMomentsMQ}
\frac{1}{\pi^*(Q)} \sum_{Q\leq p \leq 2Q} \prod_{i =1}^k f_{p}(t_i)^{n_i}= M_{\X}(\boldsymbol{n}) + O_{\boldsymbol{n}, A}\Bigl(\frac{1}{(\log Q)^{A}}\Bigr),
\end{equation}
where $M_{\X}(\boldsymbol{n})$ is defined in \eqref{eq: limiting moment}, and the implicit constant in the error term is not effective.
\end{pro}

Conditionally on GRH, the first author established this result in her thesis \cite{Hu}, with a much stronger error term $Q^{-1/2+o(1)}$. In order to compute these moments unconditionally, the main ingredient is a recent work of the second author \cite{La}, which allows us to substantially reduce the length of the Pólya Fourier expansion of $f_p$, for almost all primes $p$ such that $Q\leq p\leq 2Q$ (see Lemma \ref{lem: Tail2} below).
Another difficulty in proving Proposition \ref{thm: moments theorem} (which explains the weak error term of \eqref{AsymptoticMomentsMQ}) arises from the possible existence of Landau-Siegel zeros. Let $x$ be a large real number. By the results of~\hbox{\cite[Chap.\,20]{Da}}, it follows that for all square-free integers $|d|\leq \exp(\sqrt{\log x})$ with at most one exception $q_1$, we have
\begin{equation}\label{NoExceptional}
\sum_{p\leq x} \chi_d(p) \ll x\exp\Bigl(-c\sqrt{\log x}\Bigr),
\end{equation}
for some positive constant $c$, where $\chi_d$ denotes the Kronecker symbol attached to $d$.
Moreover, by Siegel's theorem this exceptional discriminant if it exists must satisfy
\begin{equation}\label{LowerBoundSiegel}
|q_1|\gg_A (\log x)^{4A},
\end{equation}
for any fixed constant $A>0$, where the implicit constant is not effective.
To lighten our notation we shall denote by $\chi_p:=\legendre{\cdot}{p}$ the Legendre symbol modulo a prime $p$, throughout the remaining part of the paper.
The following lemma is standard (see~\cite{Da}) but we include its proof for the sake of completeness.
\begin{lem}\label{Siegel}
Let $Q$ be large and $n\leq \exp\left(\sqrt{\log Q}\right)$ be a positive integer. Then, we~have
\begin{multline*}
\sum_{\substack{Q\leq p\leq 2Q\\ p\equiv 1 \bmod 4}} \chi_{p}(n)\\
=\begin{cases} \displaystyle{\frac{\pi^*(Q)}{2}}+O\Bigl(Q\exp\bigl(-c\sqrt{\log Q}\bigr)\Bigr), & \textup{if } n= m^2,\\
O\Bigl(Q\exp\bigl(-c\sqrt{\log Q}\bigr)\Bigr), & \textup{if } n \text{ is not of the form } m^2 \text{ or } |q_1| \cdot m^2,\end{cases}
\end{multline*}
for some positive constant $c$.
Moreover, the same estimate holds for $ \sum_{\substack{Q\leq p\leq 2Q\\ p\equiv 3 \bmod 4}} \chi_{p}(n)$.
\end{lem}
\begin{proof} We only prove the estimate for the character sum over primes $p\equiv 1 \pmod 4$ since the proof in the case $p\equiv 3 \pmod 4$ is similar.
The first estimate when $n$ is a square follows simply from the prime number theorem in arithmetic progressions together with the trivial bound $\omega(n)\leq \log n$, where $\omega(n)$ is the number of prime factors of $n$. Now, suppose that $n$ is not a square, and write $n=dm^2$ where $d$ is square-free and $d\neq |q_1|$. By the law of quadratic reciprocity we deduce that
\[
\sum_{\substack{Q\leq p\leq 2Q\\ p\equiv 1 \bmod 4}} \chi_p(n)= \sum_{\substack{Q\leq p\leq 2Q\\ p\equiv 1 \bmod 4}}\chi_p(d) + O(\omega(n))= \sum_{\substack{Q\leq p\leq 2Q\\ p\equiv 1 \bmod 4}}\chi_d(p) + O(\sqrt{\log Q}).
\]
The second estimate follows from \eqref{NoExceptional}, together with the fact that
\[\sum_{\substack{Q\leq p\leq 2Q\\ p\equiv 1 \bmod 4}}\chi_d(p)= \frac{1}{2} \sum_{Q\leq p\leq 2Q}\chi_d(p) + \frac{1}{2}\sum_{p\leq Q}\chi_{-d}(p).\qedhere
\]
\end{proof}

In order to compute the joint moments of $f_{p}(t_1), \dots, f_{p}(t_k)$ we need to considerably shorten the sum in the main term of \eqref{eq: definition Legendre path}, for all $t\in \{t_1, t_2, \dots, t_k\}$. To this end we shall use the following result from a recent work of the second author \cite{La} on the distribution of large values of quadratic character sums.

\begin{lem}\label{lem: Tail2} Let $B>0$ be a fixed constant. For all real numbers $N_1, N_2$ verifying $ \exp\big((\log\log Q)^2\big)\leq N_1 < N_2\leq 2N_1\leq Q^{21/40}$, the number of primes $p$ such that $Q\leq p\leq 2Q$ and
\[
\max_{\alpha\in [0, 1)}\biggl|\sum_{N_1\leq n\leq N_2} \frac{\chi_p(n) e(\alpha n)}{n}\biggr|\geq \frac{1}{(\log N_1)^B},
\]
is
\[
\ll_B Q\exp\Bigl(-\frac{\log Q}{10 \log\log Q}\Bigr).\]
\end{lem}
\begin{proof} This is a direct consequence of \cite[Prop.\,4.2 \& 4.3]{La} by embedding the set of primes $p$ such that $Q\leq p\leq 2Q$ into the set of all fundamental discriminants $|d|\leq 2Q$.
\end{proof}

We now have all the ingredients to prove Proposition \ref{thm: moments theorem}.

\begin{proof}[Proof of Proposition \ref{thm: moments theorem}]
Let $k$ be a positive integer, $\boldsymbol{n} = (n_1, \dots, n_k)\in (\mathbb{Z}_{\geq 0})^k$ and $0 \leq t_1 < \cdots < t_k \leq 1$ be real numbers.
Throughout the proof we shall assume that $Q$ is sufficiently large in terms of $n=n_1+\cdots+n_k$. To shorten our notation we define
\[
M_{Q}(\boldsymbol{n}):= \frac{1}{\pi^*(Q)} \sum_{Q\leq p \leq 2Q} \prod_{i =1}^k f_{p}(t_i)^{n_i}.
\]
Then, we note that
\[M_{Q}(\boldsymbol{n})= \frac{1}{2}\left(M_{Q,+}(\boldsymbol{n})+ M_{Q,-}(\boldsymbol{n})\right),\]
where
\begin{equation}\label{eq: moment sequence}
\begin{aligned}
M_{Q,+}(\boldsymbol{n})&:= \frac{2}{\pi^*(Q)} \sum_{\substack{Q\leq p \leq 2Q\\ p \equiv 1 \, \bmod 4}} \prod_{i =1}^k f_{p}(t_i)^{n_i},\\
M_{Q,-}(\boldsymbol{n})&:= \frac{2}{\pi^*(Q)} \sum_{\substack{Q\leq p \leq 2Q\\ p \equiv 3 \, \bmod 4}} \prod_{i =1}^k f_{p}(t_i)^{n_i}.
\end{aligned}
\end{equation}
Therefore, in view of \eqref{eq: MomentRandomPM} it suffices to establish the following asymptotic formulas
\begin{equation}\label{eq: AsymptoticPlusMinus}
\begin{aligned}
M_{Q, +}(\boldsymbol{n}) &= M_{\X, +}(\boldsymbol{n}) + O_{\boldsymbol{n},A}\Bigl(\frac{1}{(\log Q)^{A}}\Bigr),\\
M_{Q, -}(\boldsymbol{n}) &= M_{\X, -}(\boldsymbol{n}) + O_{\boldsymbol{n},A}\Bigl(\frac{1}{(\log Q)^{A}}\Bigr).
\end{aligned}
\end{equation}
We shall only prove this estimate for $M_{Q, +}(\boldsymbol{n})$, since the proof for $M_{Q, -}(\un)$ is similar.

Using \eqref{eq: definition Legendre path} with $Z:= Q^{21/40}$ we get
\[
M_{Q, +}(\un)= \frac{1}{\pi^n} \frac{2}{\pi^*(Q)} \sum_{\substack{Q\leq p\leq 2Q\\ p\equiv 1\bmod 4}}\prod_{i=1}^k \biggl(\sum_{a\leq Z} \chi_p(a) \frac{\sin(2\pi at_i)}{a}+ O\bigl(Q^{-1/50}\bigr)\biggr)^{n_i}.
\]
We now define $Y:=\exp((\log Q)^{1/3})$.
Let $L_1:=\lfloor \log Y /\log 2\rfloor$, $L_2:= \lfloor \log Z/\log 2\rfloor$, and put $s_{L_1}:= Y$, $s_{L_2+1}:= Z$, and $s_{\ell}:=2^{\ell}$ for $L_1+1\leq \ell\leq L_2$.
Then we have
\begin{equation}\label{PolyaSplit}
\begin{aligned}
\max_{t\in [0, 1)} \biggl|\sum_{Y\leq a\leq Z} \chi_p(a) \frac{\sin(2\pi at)}{a}\biggr|
& \leq \sum_{L_1\leq \ell \leq L_2} \max_{t \in [0, 1)} \biggl|\sum_{s_{\ell}\leq a\leq s_{\ell+1}} \frac{\chi_p(a)e(at)}{a}\biggr|.
\end{aligned}
\end{equation}
Furthermore, it follows from Lemma \ref{lem: Tail2} with the choice $B=6An+1$ that
\begin{equation}\label{BoundTailLarge}
\sum_{L_1\leq \ell \leq L_2} \max_{t \in [0, 1)} \biggl|\sum_{s_{\ell}\leq a\leq s_{\ell+1}} \frac{\chi_p(a)e(at)}{a}\biggr|
\ll
\sum_{L_1\leq \ell \leq L_2} \frac{1}{\ell^B} \ll \frac{1}{(\log Q)^{2An}},
\end{equation}
for all primes $p$ such that $Q\leq p\leq 2Q$ except for a set $\mathcal{E}(Q)$ of size
\[
|\mathcal{E}(Q)| \ll_n L_2 Q \exp\Bigl(-\frac{\log Q }{10 \log\log Q}\Bigr) \ll_n Q \exp\Bigl(-\frac{\log Q}{20 \log\log Q}\Bigr).\]
Combining the estimates \eqref{PolyaSplit} and \eqref{BoundTailLarge} we deduce that $M_{Q, +}(\un)$ equals
\begin{align*}
&\frac{2}{\pi^n\pi^*(Q)} \sum_{\substack{p\in [Q, 2Q]\setminus \mathcal{E}(Q)\\ p\equiv 1\bmod 4}}\prod_{i=1}^k \biggl(\sum_{a\leq Z} \chi_p(a) \frac{\sin(2\pi at_i)}{a}+ O(Q^{-1/50})\biggl)^{n_i}\\[-10pt]
& \hspace*{7cm} + O_n\Bigl(\exp\Bigl(-\frac{\log Q}{30\log\log Q}\Bigr)\Bigr)\\
&= \frac{2}{\pi^n\pi^*(Q)} \sum_{\substack{p\in [Q, 2Q]\setminus \mathcal{E}(Q)\\ p\equiv 1\bmod 4}}\prod_{i=1}^k \biggl(\sum_{a\leq Y} \chi_p(a) \frac{\sin(2\pi at_i)}{a}+ O\Bigl(\frac{1}{(\log Q)^{2An}}\Bigr)\biggr)^{n_i}\\[-8pt]
& \hspace*{7cm} + O_n\Bigl(\exp\Bigl(-\frac{\log Q}{30\log\log Q}\Bigr)\Bigr)\\
&= \frac{2}{\pi^n\pi^*(Q)}\sum_{\substack{p\leq Q\\ p\equiv 1\bmod 4}}\prod_{i=1}^k \biggl(\sum_{a\leq Y} \chi_p(a) \frac{\sin(2\pi at_i)}{a}+ O\Bigl(\frac{1}{(\log Q)^{2An}}\Bigr)\biggr)^{n_i}\\[-10pt]
& \hspace*{7cm} + O_n\Bigl(\exp\Bigl(-\frac{\log Q}{30\log\log Q}\Bigr)\Bigr).
\end{align*}
Expanding the main term we obtain
\begin{equation}\label{MQMain}
\begin{aligned}
M_{Q, +}(\un)
&= \frac{2}{\pi^n\pi^*(Q)}\sum_{\substack{Q\leq p\leq 2Q\\ p\equiv 1\bmod 4}}\sum_{a\leq Y^n}\chi_p(a) \frac{\B_{Y,\un, \ut, +}(a)}{a} +O _n\Bigl(\frac{(\log Y)^n}{(\log Q)^{2An}}\Bigr)\\
&= \frac{2}{\pi^n}\sum_{a\leq Y^n}\frac{\B_{Y,\un, \ut, +}(a)}{a} \frac{1}{\pi^*(Q)}\sum_{\substack{Q\leq p\leq 2Q\\ p\equiv 1\bmod 4}}\chi_p(a) +O_n\Bigl(\frac{1}{(\log Q)^{An}}\Bigr),
\end{aligned}
\end{equation}
where $\B_{Y,\un, \ut, +}(a)$ is defined by
\[\B_{Y,\un, \ut, +}(a)=
\sum_{a_{1}\cdots a_{k} = a} \prod_{i = 1}^k \sum_{\substack{b_{i,1}, \cdots b_{i,n_i} = a_i \\ b_{i,j} \leq Y}} \prod_{j = 1}^{n_i} \sin(2 \pi b_{i,j} t_i).\]
We shall first estimate the contribution of the squares $a=m^2$ which gives rise to the main term of $M_{Q, +}(\un)$. Indeed, the contribution of these terms to the right hand side of \eqref{MQMain} equals
\begin{multline*}
\frac{1}{\pi^n}\sum_{m\leq Y^{n/2}}\frac{\B_{Y,\un, \ut, +}(m^2)}{m^2} \left(1+O\left(\exp\bigl(-c_0\sqrt{\log Q}\bigr)\right)\right) \\
+ O\biggl(\frac{1}{\pi^*(Q)}\sum_{m\leq Y^{n/2}}\frac{|\B_{Y,\un, \ut, +}(m^2)|\omega(m)}{m^2}\biggr),
\end{multline*}
for some positive constant $c_0$, by the prime number theorem for arithmetic progressions.
Let $\varepsilon>0$ be small. By \eqref{eq: BoundBnt+} we have $\B_{Y,\un, \ut, +}(m^2), \B_{\un, \ut, +}(m^2)\ll_{\varepsilon, n} m^{\varepsilon}$. Since $\omega(m)\leq \log m$ and $\B_{Y,\un, \ut, +}(\ell)= \B_{\un, \ut, +}(\ell)$ if $\ell\leq Y$, we deduce that this contribution equals
\begin{equation}\label{SquaresContrib}
\begin{aligned}
\frac{1}{\pi^n}&\sum_{m\leq \sqrt{Y}}\frac{\B_{Y,\un, \ut, +}(m^2)}{m^2} + O_n\biggl(\sum_{\sqrt{Y}<m\leq Y^{n/2}}\frac{1}{m^{2-\varepsilon}}+ \exp\bigl(-c_0\sqrt{\log Q}\bigr)\biggr)\\
&= \frac{1}{\pi^n}\sum_{m=1}^{\infty}\frac{\B_{\un, \ut, +}(m^2)}{m^2} + O_n\biggl(\sum_{m> \sqrt{Y}}\frac{1}{m^{2-\varepsilon}}+\exp\bigl(-c_0\sqrt{\log Q}\bigr)\biggr)\\
&= M_{\X, +}(\un)+ O\bigl(Y^{-1/3}\bigr),
\end{aligned}
\end{equation}
by \eqref{eq: FormulaMomentRandom+}.
Next, we bound the contribution of the terms $a$ which are not of the form~$m^2$ or $|q_1|m^2$. Since $Y^n\leq \exp\left(\sqrt{\log Q}\right)$ by our assumptions on $Y$ and $Q$, it follows from Lemma \ref{Siegel} that the contribution of these terms to the main term on the right hand side of \eqref{MQMain} is
\begin{equation}\label{NonSquares}
\begin{aligned}
&\ll \exp\left(-\frac{c}{2}\sqrt{\log Q}\right)\sum_{a\leq Y^n}\frac{|\B_{Y,\un, \ut, +}(a)|}{a}\\
&\ll Y^{n/2} \exp\left(-\frac{c}{2}\sqrt{\log Q}\right) \ll \exp\Bigl(-\frac{c}{3}\sqrt{\log Q}\Bigr),
\end{aligned}
\end{equation}
since $\B_{Y,\un, \ut, +}(a)\ll \sqrt{a}$, and where $c>0$ is the constant in Lemma \ref{Siegel}.

Finally, we bound the contribution of the terms $a=|q_1|m^2$. Using the bound $\B_{Y,\un, \ut, +}(a)\ll_{n} a^{1/4}$, and bounding the inner character sum in the right hand side of \eqref{MQMain} trivially, we deduce that the contribution of these terms is
\begin{equation}\label{ExceptionalContrib}
\ll \sum_{m\leq Y^{n/2}/\sqrt{|q_1|}}\frac{\B_{Y,\un, \ut, +}(|q_1|m^2)}{|q_1|m^2}\ll_{n} |q_1|^{-3/4}\sum_{m\geq 1}\frac{1}{m^{3/2}}\ll_{A, n} (\log Q)^{-A},
\end{equation}
by \eqref{LowerBoundSiegel}.
Inserting the estimates \eqref{SquaresContrib}, \eqref{NonSquares}, and \eqref{ExceptionalContrib} in \eqref{MQMain} completes the proof.
\end{proof}
\begin{thebibliography}{BGR23}
\bibitem[10]{Ha1} A. J. Harper – “Moments of random multiplicative functions, II: High moments”, Algebra Number Theory 13 (2019), no. 10, p. 2277–2321.
\bibitem[11]{Ha2} A. J. Harper, “Moments of random multiplicative functions, I: Low moments, better than squareroot cancellation, and critical multiplicative chaos”, Forum Math. Pi 8 (2020), article no. e1 (95 pages).
\bibitem[26]{MoVaBook} H. L. Montgomery \& R. C. Vaughan, Multiplicative number theory. I. Classical theory, Cambridge Studies in Advanced Math., vol. 97, Cambridge University Press, Cambridge, 2007.
\bibitem[15]{Ka} A. Kalmynin – “Long nonnegative sums of Legendre symbols”, 2019, 23 p., arXiv:1911.10634.
\bibitem[14]{Hu} A. Hussain – “The limiting distribution of character sums”, Internat. Math. Res. Notices (2022), no. 20, p. 16292–16326.
\bibitem[19]{La} Y. Lamzouri, “The distribution of large quadratic character sums and applications”, 2022, to appear in Algebra Number Theory, arXiv:2212.03227.
\bibitem[6]{Da} H. Davenport – Multiplicative number theory, third ed., Graduate Texts in Math., vol. 74, Springer-Verlag, New York, 2000.
\end{thebibliography}
\end{document}
